\chapter{Especificación de la \acs{API} \acs{REST}ful}
\label{chap:anexo-api}

Este anexo documenta los principales \emph{endpoints} de la \ac{API}
\ac{REST}ful de \dataqual, organizados por módulo (\emph{blueprint}).
Todos los \emph{endpoints} protegidos requieren un \emph{token} \ac{JWT}
válido en la cabecera \texttt{Authorization:~Bearer~<token>}.  Las respuestas
siguen el formato estandarizado
\texttt{\{success,~data,~message\}}.


%% =========================================================================
\section{Módulo de autenticación (\texttt{/api/auth})}
\label{sec:api-auth}
%% =========================================================================

\begin{definitionlist}
  \item[\texttt{POST /api/auth/register}]  Registro de un nuevo usuario.
    Requiere \texttt{username}, \texttt{email} y \texttt{password}.

  \item[\texttt{POST /api/auth/login}]  Inicio de sesión.  Devuelve
    \emph{access\_token} y \emph{refresh\_token}.

  \item[\texttt{POST /api/auth/refresh}]  Refresco del \emph{access token}
    utilizando el \emph{refresh token}.

  \item[\texttt{POST /api/auth/logout}]  Cierre de sesión.  Revoca el
    \emph{token} actual.

  \item[\texttt{GET /api/auth/profile}]  Consulta del perfil del usuario
    autenticado.

  \item[\texttt{PUT /api/auth/profile}]  Actualización del perfil del
    usuario.
\end{definitionlist}


%% =========================================================================
\section{Módulo de proyectos (\texttt{/api/projects})}
\label{sec:api-projects}
%% =========================================================================

\begin{definitionlist}
  \item[\texttt{GET /api/projects/}]  Lista de proyectos del usuario
    autenticado, con conteo de conjuntos de datos y puntuación de la última
    evaluación.  Soporta paginación.

  \item[\texttt{POST /api/projects/}]  Creación de un nuevo proyecto.
    Requiere \texttt{name} y, opcionalmente, \texttt{description}.

  \item[\texttt{GET /api/projects/<id>}]  Detalle de un proyecto.

  \item[\texttt{PUT /api/projects/<id>}]  Actualización de un proyecto
    (nombre, descripción, configuración de métricas).

  \item[\texttt{DELETE /api/projects/<id>}]  Eliminación de un proyecto y
    todos sus recursos asociados (conjuntos de datos, evaluaciones,
    \emph{issues}).

  \item[\texttt{GET /api/projects/<id>/metrics/config}]  Consulta de la
    configuración de métricas del proyecto.

  \item[\texttt{POST /api/projects/<id>/metrics/config}]  Actualización de
    la configuración de métricas del proyecto.
\end{definitionlist}


%% =========================================================================
\section{Módulo de conjuntos de datos (\texttt{/api/datasets})}
\label{sec:api-datasets}
%% =========================================================================

\begin{definitionlist}
  \item[\texttt{POST /api/projects/<id>/datasets}]  Carga de un nuevo
    conjunto de datos (formulario \emph{multipart}).  Acepta archivos en formato \ac{CSV}.

  \item[\texttt{GET /api/datasets/<id>}]  Consulta de metadatos, esquema y
    estadísticas de un conjunto de datos.

  \item[\texttt{GET /api/datasets/<id>/versions}]  Historial de versiones de
    un conjunto de datos.

  \item[\texttt{PUT /api/datasets/<id>/sensitive-columns}]  Actualización de
    las columnas marcadas como \ac{PII}.

  \item[\texttt{GET /api/datasets/<id>/download}]  Descarga del archivo
    original.
\end{definitionlist}


%% =========================================================================
\section{Módulo de métricas (\texttt{/api/metrics})}
\label{sec:api-metrics}
%% =========================================================================

\begin{definitionlist}
  \item[\texttt{GET /api/metrics/}]  Catálogo de métricas disponibles con
    sus parámetros de configuración.

  \item[\texttt{GET /api/metrics/templates}]  Lista de plantillas de métricas
    (del sistema y del usuario).

  \item[\texttt{POST /api/metrics/templates}]  Creación de una nueva
    plantilla de métricas.

  \item[\texttt{DELETE /api/metrics/templates/<id>}]  Eliminación de una
    plantilla del usuario.
\end{definitionlist}


%% =========================================================================
\section{Módulo de evaluaciones (\texttt{/api/evaluations})}
\label{sec:api-evaluations}
%% =========================================================================

\begin{definitionlist}
  \item[\texttt{POST /api/evaluations/}]  Lanzamiento de una nueva
    evaluación.  Requiere \texttt{project\_id} y \texttt{dataset\_id}.

  \item[\texttt{GET /api/evaluations/<id>}]  Consulta de los resultados
    completos de una evaluación: métricas, puntuación, \emph{issues} y
    estado del \emph{Quality Gate}.

  \item[\texttt{GET /api/evaluations/<id>/status}]  Consulta del estado y
    progreso de una evaluación en curso.

  \item[\texttt{GET /api/evaluations/<id>/issues}]  Lista de \emph{issues}
    detectados en la evaluación, con severidad, descripción, columnas
    afectadas y muestras de filas.

  \item[\texttt{GET /api/projects/<id>/evaluations}]  Historial de
    evaluaciones de un proyecto, con puntuaciones y estados de \emph{Quality
    Gate}.
\end{definitionlist}


%% =========================================================================
\section{Módulo de administración (\texttt{/api/admin})}
\label{sec:api-admin}
%% =========================================================================

\begin{definitionlist}
  \item[\texttt{GET /api/admin/health}]  Comprobación de salud del sistema.
    Verifica la conectividad con PostgreSQL, MinIO y Redis.

  \item[\texttt{GET /api/admin/performance}]  Estadísticas de rendimiento por
    \emph{endpoint}: número de peticiones, tiempo medio, mínimo y máximo.
\end{definitionlist}


% Local Variables:
%  coding: utf-8
%  mode: latex
%  mode: flyspell
%  ispell-local-dictionary: "castellano8"
% End:
